\documentclass[10pt,a4paper]{amsart}
\usepackage[margin=19mm]{geometry}
\usepackage{amsmath,amssymb,mathtools}
\usepackage[scr=rsfs,cal=euler]{mathalfa}
\usepackage{xfrac}
\usepackage[unicode,hidelinks,bookmarks=false]{hyperref}
\newcommand{\ie}{i.e.\ }
\newcommand{\cf}{cf.\ }
\newcommand{\resp}{respectively\ }
\newcommand{\Subsection}{\subsection}
\providecommand{\nobreakdash}{\nobreak\hbox{-}}
\setlength{\emergencystretch}{2em}
\multlinegap0pt
\usepackage{bm}
\usepackage{bbm}
\usepackage{dsfont}
\usepackage{xspace}
\usepackage{cancel}
\usepackage{mathtools}
\usepackage{amsthm}
\makeatletter
\@ifundefined{definition}{\newtheorem{definition}{Definition}}{}
\@ifundefined{proposition}{\newtheorem{proposition}[definition]{Proposition}}{}
\@ifundefined{theorem}{\newtheorem{theorem}[definition]{Theorem}}{}
\makeatother
\def\Z{{\mathbb Z}}
\newcommand{\norm}[1]{\left\lVert #1 \right\rVert}
\title[Gabor Frames of Totally Positive Functions: A Complete Chara]{Gabor Frames of Totally Positive Functions: A Complete Characterization}
\author{Jaume de Dios Pont, Karlheinz Gr\"ochenig, Lukas Liehr, Irina Shafkulovska, Mitchell A. Taylor}
\date{Source: arXiv:2608.04992v1 (CC BY 4.0)}
\begin{document}
\maketitle

\noindent\emph{Source and licence.} Gabor Frames of Totally Positive Functions: A Complete Characterization. Version arXiv:2608.04992v1 (CC BY 4.0), distributed under the Creative Commons Attribution 4.0 International licence, as recorded in the arXiv licence metadata for this exact version and shown on the abstract page https://arxiv.org/abs/2608.04992v1. Full rights notice: licenses/arxiv-2608.04992-CC-BY-4.0.txt.\par\smallskip

\noindent\textit{Editorial scope.} Counted unit: the Theorem whose source label is \texttt{thm:surjectivity\_implies\_Fredholm} in the article (source0.tex lines 705--708, source label \texttt{thm:surjectivity\_implies\_Fredholm}), delivered together with the article's own complete original proof of it (source0.tex lines 710--809, one \texttt{proof} environment of 100 lines). No mathematics is written, repaired or re-derived: statement and proof are the article's own text line for line. Required context retained uncounted: eq:op\_norm\_bound (lines 555--558); prop:surjective\_band\_kernel\_bound (lines 606--612). Every definition, display and label of those lines is retained unchanged. Omitted from this package and not redistributed: 1--554, 559--605, 613--704, 709--709, 810--1819. Nothing inside a retained excerpt is omitted or altered: every retained body is delivered exactly as the article's lines are written, so no omission receipt of any kind accompanies this package. Citations: the retained excerpts carry no citation command at all, so no bibliography entry of the article is transcribed and no reference is redistributed. Reference adaptation: none was needed and none was made; every reference inside the delivered unit resolves inside it, so no unresolved reference or citation is produced. Printed slips kept literal: no printed word, symbol, hypothesis or number of the counted statement or of its proof was altered. Layout and class: the article's own class is \texttt{amsart} with packages including geometry, fontspec, minted, amsmath, amsthm, amssymb, amscd, accents, bm, amsfonts, mathtools, url, mathrsfs, dsfont; the wrapper is the lane's pinned \texttt{amsart} editorial wrapper at 10pt on A4, which restates the theorem-like environments the retained text needs and adds the article's own macro definitions that the retained text needs (\texttt{\textbackslash Z}, \texttt{\textbackslash norm}), with extra wrapper packages (bm, bbm, dsfont, xspace, cancel, mathtools). The delivered numbering is the unit's own. Assets: no retained span contains a figure environment, a graphics-inclusion command, a PDF-page inclusion, a table or a TikZ picture; no image file is copied into this package, the package declares no asset dependency and no image file is redistributed with it. Licence: version arXiv:2608.04992v1 of this article is distributed under the Creative Commons Attribution 4.0 International licence, as recorded in the arXiv licence metadata for this exact version and shown on the abstract page https://arxiv.org/abs/2608.04992v1. A full rights notice is delivered at licenses/arxiv-2608.04992-CC-BY-4.0.txt. Characters: every retained excerpt is pure ASCII, so the delivered engine, XeLaTeX with its default Latin Modern text font, reports no missing character.
\begin{equation}\label{eq:op_norm_bound}
    \| A \|_{1 \to 1} 
    =\sup_{l\in\mathbb Z}\sum_{k\in\mathbb Z}|A_{kl}|.
\end{equation}

\begin{proposition}\label{prop:surjective_band_kernel_bound}
    Every surjective band operator $A\in\mathcal L(\ell^1(\Z))$ of
bandwidth $w\in\mathbb N_0$ satisfies
$
\dim\ker A\leq 2w .
$
\end{proposition}

\begin{theorem}\label{thm:surjectivity_implies_Fredholm}
If $A = (A_{kl})_{k,l \in \Z}$ possesses  polynomial off-diagonal
decay $\sigma >1$ and  is surjective on $\ell^1(\Z)$, then $A$ is Fredholm on $\ell^1(\Z)$. In particular, $A$ has a finite-dimensional kernel. 
\end{theorem}

\begin{proof}
It suffices to show that $A$ has a finite-dimensional kernel.

\textbf{Step 1. A perturbation argument.}
Let $\tilde{A},R\in \mathcal{L}(\ell^1(\Z))$ with
\begin{equation} \label{eq:c1}
    AR=\mathrm{Id}_{\ell^1(\Z)},\quad \|A-\tilde{A}\|_{1\to 1} \norm{R}_{1\to 1} <1.
\end{equation}
We claim that there exists a bounded invertible operator $\Phi$ with $A=\tilde{A}\Phi$. In particular, $\tilde{A}$ is surjective and
\begin{equation}\label{eq:same_kernel}
    \dim \ker A= \dim\ker \tilde{A} .
\end{equation}

Set $E=A-\tilde{A}$. By assumption~\eqref{eq:c1}, $\|ER\|_{1\to 1}<1$. 
Thus, the operator $\mathrm{Id}_{\ell^1(\Z)}- ER$ is invertible. Hence, from 
\begin{equation}
    \tilde{A}R = (A- E)R =\mathrm{Id}_{\ell^1(\Z)} -  ER,
\end{equation}
it follows that $\tilde{A}$ is surjective.
We set 
\begin{align}
    \Phi & = \mathrm{Id}_{\ell^1(\Z)} +R\sum\limits_{l=0}^\infty (ER)^l E
     = \mathrm{Id}_{\ell^1(\Z)} +\sum\limits_{l=1}^\infty (RE)^l = (\mathrm{Id}_{\ell^1(\Z)} -RE)^{-1}.
\end{align}
Clearly, $\Phi$ is invertible with inverse $\mathrm{Id}_{\ell^1(\Z)}
-RE$, and thus 
\begin{equation}
    A\Phi^{-1} = A(\mathrm{Id}_{\ell^1(\Z)} -RE) = A- ARE = A-E = \tilde{A}.
\end{equation}
% \newpage
\textbf{Step 2. Existence of a right inverse. }
Since $A$ is surjective, the open mapping theorem implies that there exists a constant
$C>0$ such that for every $c\in {\ell^1(\Z)}$ there exists $b\in {\ell^1(\Z)}$ with
$$
Ab=c,
\quad
\|b\|_1\leq C\|c\|_1 .
$$
For each $l\in\mathbb Z$, choose $b^{(l)}\in {\ell^1(\Z)}$ such that
$$
Ab^{(l)}=e_l,
\quad
\|b^{(l)}\|_1\leq C,
$$
and define
$$
Rc:=\sum_{l\in\mathbb Z}c_l b^{(l)} .
$$
This series converges absolutely in ${\ell^1(\Z)}$, and
$$
\|Rc\|_1\leq C\|c\|_1 .
$$
Thus, $R\in\mathcal L({\ell^1(\Z)})$. Since the defining series is absolutely
convergent, applying $A$ termwise gives
$$
ARc=\sum_{l\in\mathbb Z}c_lAb^{(l)}
=\sum_{l\in\mathbb Z}c_le_l
=c.
$$
Therefore, we have $AR=\mathrm{Id}_{\ell^1(\Z)}$.

\textbf{Step 3. Approximation by band operators.}
Let $A_n$ be the band truncation of $A$  defined by
$$
(A_n)_{kl}=
\begin{cases}
A_{kl}, & |k-l|\leq n,\\
0, & |k-l|>n .
\end{cases}
$$
Then 
$
\|A-A_n\|_{1\to1} \to 0
$
as $n \to \infty$. Indeed, 
   let $C>0$, $\sigma>1$ be such that 
   \begin{equation}
       |A_{kl}|\le C(1+|k-l|)^{-\sigma},
\quad k,l\in\mathbb Z.
   \end{equation}
   Then by \eqref{eq:op_norm_bound}, we have
   \begin{align}
       \norm{A-A_n}_{1 \to 1} & =\sup\limits_{l\in\Z}\sum\limits_{k\in\Z}|(A-A_n)_{kl}| = \sup\limits_{l\in\Z}\sum\limits_{|k-l|>n}|A_{kl}| \\
       & \leq  C \sup\limits_{l\in\Z}\sum\limits_{|k-l|>n} (1+|k-l|)^{-\sigma} 
       = C \sum\limits_{|j|>n} (1+|j|)^{-\sigma}.
   \end{align}
   Since $\sum\limits_{j\in\Z} (1+|j|)^{-\sigma}$ converges, $\norm{A-A_n}_{1\to 1}\to 0$ as $n\to \infty$.

\textbf{Step 4.}
Fix a large enough $n$ so that
$$
\|A-A_n\|_{1\to1}\,\|R\|_{1\to1}<1 .
$$

By Step 1, $A_n$ is surjective. 
Applying \eqref{eq:same_kernel} and Proposition~\ref{prop:surjective_band_kernel_bound} to the surjective band operator $A_n$, we obtain
\begin{equation}
    \dim\ker A= \dim\ker A_n \leq 2n.
\end{equation}
\end{proof}

\end{document}
